\documentclass[10pt,a4paper]{amsart}
\usepackage{amsmath,amssymb,amsthm,graphicx}
\usepackage{verbatim}
\newenvironment{reference}{\comment}{\endcomment}
%\newenvironment{reference}{}{}
\newenvironment{slogan}{\comment}{\endcomment}
\newenvironment{history}{\comment}{\endcomment}

% For commutative diagrams we use Xy-pic
\usepackage[all]{xy}

% We use 2cell for 2-commutative diagrams.
\xyoption{2cell}
% Theorem environments.
%
\theoremstyle{plain}
\newtheorem{theorem}[subsection]{Theorem}
\newtheorem{proposition}[subsection]{Proposition}
\newtheorem{lemma}[subsection]{Lemma}

\theoremstyle{definition}
\newtheorem{definition}[subsection]{Definition}
\newtheorem{example}[subsection]{Example}
\newtheorem{exercise}[subsection]{Exercise}
\newtheorem{situation}[subsection]{Situation}

\theoremstyle{remark}
\newtheorem{remark}[subsection]{Remark}
\newtheorem{remarks}[subsection]{Remarks}

\numberwithin{equation}{subsection}

% Macros
%
\def\lim{\mathop{\mathrm{lim}}\nolimits}
\def\colim{\mathop{\mathrm{colim}}\nolimits}
\def\Spec{\mathop{\mathrm{Spec}}}
\def\Hom{\mathop{\mathrm{Hom}}\nolimits}
\def\Ext{\mathop{\mathrm{Ext}}\nolimits}
\def\SheafHom{\mathop{\mathcal{H}\!\mathit{om}}\nolimits}
\def\SheafExt{\mathop{\mathcal{E}\!\mathit{xt}}\nolimits}
\def\Sch{\mathit{Sch}}
\def\Mor{\mathop{\mathrm{Mor}}\nolimits}
\def\Ob{\mathop{\mathrm{Ob}}\nolimits}
\def\Sh{\mathop{\mathit{Sh}}\nolimits}
\def\NL{\mathop{N\!L}\nolimits}
\def\CH{\mathop{\mathrm{CH}}\nolimits}
\def\proetale{{pro\text{-}\acute{e}tale}}
\def\etale{{\acute{e}tale}}
\def\QCoh{\mathit{QCoh}}
\def\Ker{\mathop{\mathrm{Ker}}}
\def\Im{\mathop{\mathrm{Im}}}
\def\Coker{\mathop{\mathrm{Coker}}}
\def\Coim{\mathop{\mathrm{Coim}}}

% Boxtimes
%
\DeclareMathSymbol{\boxtimes}{\mathbin}{AMSa}{"02}

%
% Macros for moduli stacks/spaces
%
\def\QCohstack{\mathcal{QC}\!\mathit{oh}}
\def\Cohstack{\mathcal{C}\!\mathit{oh}}
\def\Spacesstack{\mathcal{S}\!\mathit{paces}}
\def\Quotfunctor{\mathrm{Quot}}
\def\Hilbfunctor{\mathrm{Hilb}}
\def\Curvesstack{\mathcal{C}\!\mathit{urves}}
\def\Polarizedstack{\mathcal{P}\!\mathit{olarized}}
\def\Complexesstack{\mathcal{C}\!\mathit{omplexes}}
% \Pic is the operator that assigns to X its picard group, usage \Pic(X)
% \Picardstack_{X/B} denotes the Picard stack of X over B
% \Picardfunctor_{X/B} denotes the Picard functor of X over B
\def\Pic{\mathop{\mathrm{Pic}}\nolimits}
\def\Picardstack{\mathcal{P}\!\mathit{ic}}
\def\Picardfunctor{\mathrm{Pic}}
\def\Deformationcategory{\mathcal{D}\!\mathit{ef}}
\usepackage{cleveref}
\crefname{originallemma}{original lemma}{original lemmas}
\Crefname{originallemma}{Original lemma}{Original lemmas}
\makeatletter
\newlabel{categories-lemma-yoneda}{{001P}{source}{}{sourcepointer-1}{}}
\newlabel{categories-lemma-yoneda@cref}{{[originallemma][0][]001P}{[1][0][]source}{}{}{}}
\makeatother
\title{A left adjoint defined on a generating set of quotients is everywhere defined}
\author{The Stacks Project Authors}
\date{stacks-project:homology.tex:local-snapshot-20261012; GNU FDL 1.2 or later}
\begin{document}
\maketitle
\noindent\textit{Editorial note.} All delivered mathematical text is the author's own, copied from the pinned Stacks Project chapter file homology.tex; no statement or proof step was rewritten, reordered or removed.\par\medskip
\noindent\textit{Editorial note.} The proof is the author's complete argument; every step is written out inside the counted span, the closing sentence 'Hence we see that we may take P = Ob(B), i.e., we see that v is everywhere defined.' is the conclusion of the argument rather than a placeholder, and no step is passed to the reader.\par\medskip
\noindent\textit{Editorial note.} Every labelled reference occurring inside the counted spans is carried in assembly.external\_references as a non-load-bearing navigation pointer (same-chapter ones at their definition site, cross-chapter ones at the invocation inside the counted span), so no load-bearing mathematics is replaced by a stub. No bibliography citation occurs inside the counted spans.\par\medskip
\noindent\textit{Editorial note.} The same-name check for this card scanned every staged Stacks chapter of this repository -- all sources/stacks-*/pinned-original-source.tex files, the card's own chapter excluded -- and the labels found are recorded in dedup.same\_name\_cross\_chapter on the card.\par\medskip
\noindent\textit{Editorial note.} The preamble is the chapter's own native preamble excerpt already used for this source: the comment-type reference, slogan and history environments (native preamble lines 8-13), the Xy-pic setup (native preamble lines 14-18) and the native theorem declarations with the mathematical macros (native preamble lines 154-217, the actual last line of the file).\par\medskip
\noindent\textit{Allowed standard prerequisites (editorial).}\begin{itemize}
\item The Yoneda lemma. Categories, Lemma categories-lemma-yoneda. Identifies the object v(P) representing the prescribed functor A -> Hom\_B(P, u(A)) up to unique isomorphism, which is how the functor is defined on the chosen set of objects.
\item Quotients and exact sequences in an abelian category. Same chapter, the standard description of an object as the quotient of a map onto it and the exactness of the resulting sequence. Lets the author choose an exact sequence P\_2 -> P\_1 -> B -> 0 for an arbitrary object B, which is what extends the construction beyond the chosen set.
\end{itemize}\medskip
\ifdefined\hypertarget\hypertarget{sourcepointer-1}{}\fi
\noindent\textit{Source reference.} Original lemma 001P: sources/stacks-categories/pinned-original-source.tex, source lines 758--758. Navigation pointer invoked at proof line 7622; the counted span contains the pointer only and the cited result is not reproduced inside it, so the pointer is registered as a non-load-bearing navigation stub in the established form for Stacks material. No bibliography citation is involved.\par\medskip
\begin{lemma}
\label{lemma-partially-defined-adjoint}
Let $\mathcal{A}$ and $\mathcal{B}$ be abelian categories.
Let $u : \mathcal{A} \to \mathcal{B}$ be a functor.
If there exists a subset $\mathcal{P} \subset \Ob(\mathcal{B})$
such that
\begin{enumerate}
\item every object of $\mathcal{B}$ is a quotient of an element
of $\mathcal{P}$, and
\item for every $P \in \mathcal{P}$ there exists an object
$Q$ of $\mathcal{A}$ such that
$\Hom_\mathcal{A}(Q, A) = \Hom_\mathcal{B}(P, u(A))$ functorially
in $A$,
\end{enumerate}
then there exists a left adjoint $v$ of $u$.
\end{lemma}

\begin{proof}
By the Yoneda lemma (Categories, Lemma \ref{categories-lemma-yoneda})
the object $Q$ of $\mathcal{A}$ corresponding to $P$ is defined up to
unique isomorphism by the formula
$\Hom_\mathcal{A}(Q, A) = \Hom_\mathcal{B}(P, u(A))$. Let us write
$Q = v(P)$. Denote $i_P : P \to u(v(P))$ the map corresponding to
$\text{id}_{v(P)}$ in $\Hom_\mathcal{A}(v(P), v(P))$. Functoriality
in (2) implies that the bijection is given by
$$
\Hom_\mathcal{A}(v(P), A) \to \Hom_\mathcal{B}(P, u(A)),\quad
\varphi \mapsto u(\varphi) \circ i_P
$$
For any pair of elements $P_1, P_2 \in \mathcal{P}$ there is a canonical map
$$
\Hom_\mathcal{B}(P_2, P_1)
\to
\Hom_\mathcal{A}(v(P_2), v(P_1)),\quad
\varphi \mapsto v(\varphi)
$$
which is characterized by the rule
$u(v(\varphi)) \circ i_{P_2} = i_{P_1} \circ \varphi$ in
$\Hom_\mathcal{B}(P_2, u(v(P_1)))$.
Note that $\varphi \mapsto v(\varphi)$ is
compatible with composition; this can be seen directly
from the characterization. Hence $P \mapsto v(P)$ is a functor from
the full subcategory of $\mathcal{B}$ whose objects are the elements
of $\mathcal{P}$.

\medskip\noindent
Given an arbitrary object $B$ of $\mathcal{B}$ choose an exact sequence
$$
P_2 \to P_1 \to B \to 0
$$
which is possible by assumption (1). Define $v(B)$ to be the object of
$\mathcal{A}$ fitting into the exact sequence
$$
v(P_2) \to v(P_1) \to v(B) \to 0
$$
Then
\begin{align*}
\Hom_\mathcal{A}(v(B), A)
& =
\Ker(\Hom_\mathcal{A}(v(P_1), A) \to \Hom_\mathcal{A}(v(P_2), A)) \\
& =
\Ker(\Hom_\mathcal{B}(P_1, u(A)) \to \Hom_\mathcal{B}(P_2, u(A))) \\
& =
\Hom_\mathcal{B}(B, u(A))
\end{align*}
Hence we see that we may take $\mathcal{P} = \Ob(\mathcal{B})$, i.e., we
see that $v$ is everywhere defined.
\end{proof}

\end{document}
